\chapter{Methodik}\label{ch:methods}

In diesem Kapitel wird die methodische Vorgehensweise zur Untersuchung der Forschungsfrage vorgestellt. Ziel ist es, die Leistungsfähigkeit des entwickelten Reinforcement Learning-Ansatzes für die zwei Aufgabenbereiche – Navigation eines Unitree Go2 Edu in simulierten Umgebungen sowie taktile Interaktion durch Kopfberührung – systematisch zu evaluieren. Dazu werden verschiedene Trainingsläufe mit unterschiedlichen Zufallsinitialisierungen (Seeds) durchgeführt, um die Robustheit und Reproduzierbarkeit der Ergebnisse zu gewährleisten.

Die folgenden Abschnitte beschreiben zunächst die verwendete Simulationsumgebung und das Robotermodell, anschließend den eingesetzten Lernalgorithmus sowie die Gestaltung der Belohnungsfunktion und des Curriculums. Abschließend werden der konkrete Versuchsaufbau und die verwendeten Evaluationsmetriken erläutert.
\section{Die Simulationsumgebung Genesis und das Robotermodell}
Für die Konzeption, das Training und die Evaluation der DRL-Policies wird die physikbasierte Simulationsplattform \textbf{Genesis} eingesetzt \cite{authorsGenesisGenerativeUniversal2024}. Die Wahl von Genesis gegenüber traditionellen Robotik-Simulatoren (wie Gazebo oder PyBullet) begründet sich durch spezifische architektonische und leistungstechnische Merkmale:
\begin{enumerate}
	\item \textbf{GPU-beschleunigte Physik-Pipeline:} Genesis basiert auf einem hochoptimierten CUDA- und Taichi-Backend, welches die numerische Integration von Starrkörpermechanik und Kontaktdynamik vollständig auf der GPU ausführt.
	\item \textbf{Massiv-parallele Trainingsumgebungen auf Consumer-Hardware:} Die Plattform ist auf die simultane Simulation einer großen Anzahl paralleler Roboter-Instanzen auf einer einzelnen Desktop-Grafikkarte ausgelegt (der vom Framework angegebene Durchsatz liegt nach Herstellerangaben bei einer Größenordnung von hunderten bis tausenden Instanzen bzw. mehreren zehntausend Zeitschritten pro Sekunde ohne dedizierte Server-Cluster \cite{authorsGenesisGenerativeUniversal2024}). Diese Werte wurden im Rahmen dieser Arbeit nicht eigenständig durch einen Benchmark verifiziert; sie dienen lediglich der Einordnung der Plattform.
	\item \textbf{Native PyTorch-Tensor-Schnittstelle:} Beobachtungen und Aktionsvektoren verbleiben ohne rechenintensive CPU-GPU-Speichertransfers im VRAM der Grafikkarte, was den Datendurchsatz zwischen der Physik-Engine und dem neuronalen Netzwerk maximiert.
\end{enumerate}

Neben Genesis existieren weitere moderne GPU-basierte Plattformen für massiv-paralleles Reinforcement Learning, darunter \textbf{NVIDIA Isaac} (Isaac Sim/Lab, auf Omniverse aufbauend), die ebenfalls PyTorch-anbindbare Umgebungen bereitstellen \cite{mittal2025isaaclab}. Für die hier durchgeführte Machbarkeitsstudie bot sich Genesis dennoch als die geeignete Wahl an. Zum einen liegt mit der offiziellen PPO-Referenzimplementierung des Frameworks eine etablierte und dokumentierte Trainingspipeline vor, auf der das Training dieser Arbeit aufbaut \cite{authorsGenesisGenerativeUniversal2024}. Zum anderen erfüllt Genesis mit seiner GPU-beschleunigten, headless ausführbaren und quelloffenen Architektur die Anforderungen der hier vorgesehenen Versuche vollständig. Eine vergleichende Bewertung oder ein Benchmark der Plattformen wurde im Rahmen dieser Arbeit nicht durchgeführt.

Als Roboterplattform kommt das digitale Modell des \textbf{Unitree Go2 Edu} zum Einsatz \cite{unitreeGo2Unitree2023}. Der Go2 Edu ist ein hochagiler quadrupedaler Roboter mit einer Gesamtmasse von rund $15\,\text{kg}$ und 12 aktiven Freiheitsgraden (drei rotatorische Gelenke pro Bein: Hüftabduktion/Adduktion HAA, Hüftflexion/Extension HFE und Knieflexion/Extension KFE). Die Positionsregelung der Gelenke erfolgt über individuell parametrisierte PD-Regler bei einer Steuerfrequenz von $50\,\text{Hz}$ ($\Delta t = 0{,}02\,\text{s}$). 

In der Simulation erhält die Policy zu jedem Zeitschritt einen Beobachtungsvektor bestehend aus propriozeptiven Daten (Gelenkwinkel $q$, Gelenkwinkelgeschwindigkeiten $\dot{q}$, Rumpforientierung in Euler-Winkeln $[\phi, \theta, \psi]$ und Winkelgeschwindigkeiten $\vec{\omega}_{\text{base}}$), aufgabenspezifischen Hindernissignalen sowie bei der Interaktion binären Zustandsindikatoren für Berührung und aktive Geste. Die Berührung wird in der verwendeten Implementierung positionsbasiert über den Abstand einer virtuellen Handkugel zum Kopf erkannt. Kontaktkräfte werden von der Simulation zwar registriert, das für die Policy verwendete Berührungssignal basiert jedoch ausschließlich auf dieser positionsbasierten Kontakt-/Abstandserkennung; eine Kraftmessung ist nicht Bestandteil des Signals. Auf Domain Randomization wurde im Rahmen dieser Machbarkeitsstudie bewusst verzichtet.
\section{Der Reinforcement-Learning-Algorithmus}
Für das Training der Steuerungspolicies wird der PPO-Algorithmus eingesetzt, welcher bereits in \autoref{ss:PPO} erläutert wurde. Die Algorithmenwahl orientiert sich primär an der offiziellen Referenzimplementierung des Genesis-Frameworks \cite{authorsGenesisGenerativeUniversal2024}, welche PPO als Standard-Lernalgorithmus vorgibt und auf dessen Hyperparameterkonfiguration das Training in dieser Arbeit aufbaut. Die Entscheidung, dieser Vorgabe zu folgen, ist methodisch begründet: Sie vermeidet unnötige Komplexität bei der Algorithmenauswahl für eine Machbarkeitsstudie und erlaubt eine direkte Vergleichbarkeit mit dem Framework-eigenen Referenzverhalten. Darüber hinaus ist PPO in der Literatur für massiv-parallele Locomotion-Aufgaben gut dokumentiert und erprobt \cite{rudin2022learningwalkminutesusing, schwarke2025rslrllearninglibraryrobotics}.

\begin{table}[htbp]
\centering
\small
\begin{tabular}{lcc}
\toprule
\textbf{Hyperparameter} & \textbf{Navigation} & \textbf{Interaktion} \\
\midrule
Clip-Parameter $\epsilon$ & 0,20 & 0,20 \\
Lernrate & $1{,}0\cdot10^{-3}$ & $1{,}0\cdot10^{-4}$ \\
Diskontfaktor $\gamma$ & 0,99 & 0,99 \\
GAE-Parameter $\lambda$ & 0,95 & 0,95 \\
Entropiekoeffizient & 0,001--0,02$^\ast$ & 0,01 \\
Wertfunktionskoeffizient & 1,0 & 0,3 \\
Lernepochen pro Update & 8 & 5 \\
Mini-Batches pro Update & 4 & 8 \\
Schritte pro Umgebung & 48 & 32 \\
Aktivierung & ELU & ELU \\
Actor-/Critic-Schichten & 512--256--128 & 512--256--128 \\
\bottomrule
\end{tabular}
\caption{Verwendete PPO- und Netzwerkparameter. $^\ast$Der Entropiekoeffizient der Navigation wurde je nach Curriculum-Stufe angepasst: Stehen 0{,}001, Laufen 0{,}02, Ausweichen und Springen 0{,}005.}
\label{tab:ppo_hyperparameters}
\end{table}

\section{Belohnungsfunktionen} \label{reward_functions}

Die Gestaltung der Belohnungsfunktion (Reward Shaping) bildet ein zentrales Steuerungsinstrument, um das Verhalten des Agenten in hochdimensionalen Aktionsräumen zu beeinflussen. Da spärliche Belohnungssignale (Sparse Rewards) die Exploration erschweren können, wird in dieser Arbeit ein modularer, überwiegend dicht belohnter Ansatz (Dense Rewards) verfolgt.

Die Gesamtbelohnung $r_t$ setzt sich zum Zeitschritt $t$ aus einer gewichteten Summe von Teilbelohnungen $r_{i,t}$ zusammen:
\begin{equation}
r_t = \sum_{i=1}^{n} w_i \cdot r_{i,t}
\end{equation}
wobei die Gewichte $w_i \in \mathbb{R}$ den relativen Einfluss und die Priorisierung der einzelnen Zielkomponenten festlegen.

Für die \textbf{Navigationsumgebung} unterteilt sich die Belohnungsarchitektur in vier funktionale Kategorien:
\begin{itemize}
	\item \textbf{Zielannäherung und Fortschritt ($r_{\text{target\_proximity}}$, $r_{\text{target\_reached}}$, $r_{\text{target\_progress}}$, $r_{\text{goal\_heading}}$):} Belohnt die Reduktion der euklidischen Distanz zum Navigationsziel sowie das direkte Erreichen des Zielradius. Dies erzeugt den primären Vorwärtsdrang und verhindert zielloses Verharren.
	\item \textbf{Hindernisvermeidung und Trajektorienführung ($r_{\text{blocker\_clearance}}$, $r_{\text{centerline\_near\_blocker}}$):} Bestraft zu geringe Abstände zum Hindernis sowie starke seitliche Abweichungen vom direkten Pfad vor dem Passieren des Hindernisses, um ein kontrolliertes Umgehen zu erzwingen.
	\item \textbf{Haltungskontrolle und Stabilität ($r_{\text{height}}$, $r_{\text{upright}}$, $r_{\text{symmetry}}$):} Forciert eine Soll-Körperhöhe ($h_{\text{target}} = 0{,}42\,\text{m}$), minimiert Roll- und Nickwinkelabweichungen ($\phi, \theta \approx 0$) und fördert symmetrische Gelenkbewegungen zur Vermeidung von Gangasymmetrien und Stürzen.
	\item \textbf{Sprung- und Phasenbelohnungen ($r_{\text{jump\_approach}}$, $r_{\text{jump\_takeoff}}$, $r_{\text{jump\_flight}}$, $r_{\text{landing\_stability}}$):} In der Sprungstufe werden Annäherung, vertikale Bewegung, ausreichende Rumpfhöhe und Stabilität im Landeabschnitt phasenabhängig bewertet. Die Terme garantieren nicht, dass alle vier Beine gleichzeitig abheben.
\end{itemize}

Für die \textbf{Interaktionsumgebung (Streicheln/Petting)} fokussiert die Belohnungsfunktion auf die positionsbasierte Berührungserkennung und die Reaktion während des aktiven Kontakts:
\begin{itemize}
	\item \textbf{Kontakt- und Gestenreaktion ($r_{\text{gesture\_during\_touch}}$, $r_{\text{gesture\_movement}}$):} Das Kontakt-Signal aktiviert die Geste; die ausgewertete Konfiguration belohnt während der aktiven Geste den Betrag ausgewählter Gelenkgeschwindigkeiten. Ein expliziter Trajektorien- oder Abschlussreward ist nicht implementiert.
	\item \textbf{Standstabilität und Sitzvermeidung ($r_{\text{no\_sitting}}$, $r_{\text{petting\_stability}}$, $r_{\text{flexible\_height}}$):} Bestraft ein Absenken des Rumpfes auf den Boden (Sitzen/Einknicken) und stellt sicher, dass der Roboter während der Berührung aufrecht und stabil steht.
\end{itemize}

Die detaillierte mathematische Definition aller Teilterme, deren physikalische Messgrößen, die Begründung der gewählten Gewichtungsfaktoren sowie eine vergleichende Analyse von zielgerichteten Belohnungen und potenziellen Reward-Hacking-Effekten sind in \autoref{ch:reward_function_maths} dokumentiert. Der vollständige Quellcode ist im Projekt-Repository verfügbar.\footnote{Projekt-Repository mit Quellcode und Konfigurationen: \url{https://github.com/w0lkensucher/reinforcement_learning_thesis}}

\section{Netzwerk-Architektur}

Für die Implementierung der Policy- und Value-Netzwerke wurden vollständig verbundene neuronale Netzwerke verwendet. Beide Netzwerke besitzen drei Hidden Layer mit 512, 256 und 128 Neuronen, die durch ELU-Funktionen aktiviert werden. Die Ausgabeschicht ist jeweils linear ohne Aktivierungsfunktion. Die konkreten Parameter sind in \autoref{tab:ppo_hyperparameters} dokumentiert.

Die Policy gibt kontinuierliche Aktionen für die 12 Gelenke aus. Diese werden mit \texttt{action\_scale} skaliert und als Gelenkpositionsziele an die PD-Regler übergeben; direkte Drehmomentaktionen der Policy werden nicht verwendet. Für die Aktionsverteilung kommt eine Gaußsche Normalverteilung mit trainierten Varianzparametern zum Einsatz. Das Value-Netzwerk schätzt den Erwartungswert der kumulierten Belohnung.

Diese Architektur orientiert sich an den Standardimplementierungen im Genesis-Framework \cite{authorsGenesisGenerativeUniversal2024} und bietet einen ausreichenden Kompromiss zwischen Modellflexibilität und Trainingsstabilität für kontinuierliche Steuerungsaufgaben.

\section{Die Trainingsstrategie: Curriculum Learning}

Um den Lernprozess in hochdimensionalen kontinuierlichen Aktionsräumen zu strukturieren und das Problem spärlicher Belohnungssignale (Sparse Rewards) bei komplexen Bewegungsabläufen zu überwinden, wird ein \textbf{Curriculum-Learning-Ansatz} eingesetzt \cite{heessEmergenceLocomotionBehaviours2017, rudin2022learningwalkminutesusing}. Das Grundprinzip besteht darin, die Komplexität der Aufgabenstellung schrittweise zu steigern: Der Agent erlernt zunächst fundamentale Basisfähigkeiten (wie Gleichgewicht und Rumpfstabilisierung), bevor diese als funktionale Bausteine für anspruchsvollere Navigations- und dynamische Interaktionsaufgaben wiederverwendet werden.

\autoref{tab:curriculum_stages} gibt einen systematischen Überblick über die konzipierte Stufenarchitektur beider Trainingsumgebungen.

\begin{table}[htbp]
\centering
\small
\begin{tabular}{cllll}
\toprule
\textbf{Stufe} & \textbf{Bezeichnung} & \textbf{Lernziel / Kernanforderung} & \textbf{Initialisierung / Basis} & \textbf{Horizon} \\
\midrule
\multicolumn{5}{l}{\textit{Navigationsumgebung}} \\
1 & Stehen & Rumpfhöhe $0{,}42\,\text{m}$, Stabilität, $\vec{v} = 0$ & Zufällige Gewichtsinitialisierung & 400 Iter. \\
2 & Laufen & Tracking von $v_x, v_y, \omega_z$, Trabgang & Checkpoint Stufe 1 (Iter. 50) & 1049 Iter. \\
3 & Navigation & Zielansteuerung, Umgehen von Hindernis & Checkpoint Stufe 2 (Iter. 1049) & 1500 Iter. \\
3b (4) & Springen & Vertikalsprung, Flugphase, Landung & Checkpoint Stufe 2 (Iter. 1049) & 1500 Iter. \\
\midrule
\multicolumn{5}{l}{\textit{Interaktionsumgebung (Petting)}} \\
1 & Stehen & Standstabilität unter externen Kräften & Zufällige Gewichtsinitialisierung & 300 Iter. \\
2 & Gestenreaktion & Taktile Berührungserkennung \& Geste & Checkpoint Stufe 1 (Iter. 50) & 600 Iter. \\
\bottomrule
\end{tabular}
\normalsize
\caption{Übersicht der Curriculum-Stufenarchitektur für Navigation und Interaktion}
\label{tab:curriculum_stages}
\end{table}

\subsection{Stufenarchitektur und Verzweigung}
Jede Curriculum-Stufe baut methodisch auf den Gewichten der jeweils vorangegangenen Stufe auf. Der Wechsel zur nächsten Stufe erfolgte manuell, wenn die betrachteten Kurven über einen Beobachtungsabschnitt keine deutliche weitere Verbesserung zeigten. Dies wird im Folgenden als stabiles Plateau bezeichnet und nicht als formaler Konvergenznachweis.

Hierbei ist zu beachten, dass das \textbf{Springen (Stufe 4)} strukturell keine rein sequentielle Fortsetzung der horizontalen Navigation (Stufe 3) darstellt, sondern als \textbf{Verzweigung 3b} auf Grundlage der erlernten Laufbewegung (Stufe 2) operiert. Während Stufe 3 das zweidimensionale Ausweichen auf horizontaler Ebene schult, verlangt Stufe 3b den Übergang in die vertikale Dimension unter temporärer Aufgabe des stabilen Bodenkontakts.

In der Interaktionsumgebung dient Stufe 1 der Etablierung robuster Standstabilität unter dem Einfluss gravitativer und externer Störkräfte. In Stufe 2 wird der Roboter darauf trainiert, auf physische Berührungen am Kopf (simuliert durch eine taktile Kollisionsgeometrie) mit einer koordinierten Gelenkbewegung zu reagieren. Die Herausforderung besteht darin, zwischen Berührungs- und Nicht-Berührungsphasen zu differenzieren und während der aktiven Gestenausführung das Gleichgewicht nicht zu verlieren.

\section{Versuchsaufbau und Evaluationsmetrik}

Zur methodisch fundierten Bewertung der trainierten Policies wird ein strukturierter Evaluationsansatz verwendet. Trainingsläufe im Reinforcement Learning können aufgrund zufälliger Netzwerkinitialisierung und stochastischer Aktionsauswahl unterschiedliche Ergebnisse liefern \cite{hendersonDeepReinforcementLearning2019}. In den ausgewerteten Versuchsreihen wurden die Seeds 1, 8, 42, 24 und 77 verwendet. Die Grundstufen der Navigation (Stand, Laufen) sowie beide Interaktionsstufen (Stehen, Geste) wurden jeweils mit allen fünf Seeds vollständig trainiert. Für die darauf aufbauenden Stufen (Ausweichen, Springen) wurde konsistent vom zuvor ausgewählten Lauf-Checkpoint (Seed 1, Iteration 1049) ausgegangen und jeweils mit fünf Seeds weitertrainiert; die fünf Läufe dieser Folgestufen teilen sich also denselben Startzustand, unterscheiden sich jedoch in den Zufalls-Seeds des Folgetrainings. Diese Zuordnung der Checkpoints sowie die quantitativen Angaben sind eigene Versuchsdaten und werden im Ergebniskapitel dokumentiert. Die Standardkonfigurationen definieren jeweils nur einen Seed pro Trainingsprozess; aus der Konfiguration allein kann daher nicht auf fünf vollständige, unabhängige Wiederholungen jeder Stufe geschlossen werden.

Für die Einordnung dieser Versuche ist zu beachten, dass die verfügbaren Läufe, anders als bei der hier bewusst nicht eingesetzten Domain Randomization (vgl. Grundlagen in \autoref{ch:theoretical_foundation}), nur eine orientierende Betrachtung der Trainingsverläufe erlauben und keine belastbaren inferenzstatistischen Aussagen begründen.

Die Bewertung der gelernten Verhaltensweisen basiert auf zwei sich ergänzenden Ebenen: einer \textbf{qualitativen Inspektion} aufgezeichneter Episoden und den dazu korrespondierenden \textbf{quantitativen Trainingsmetriken}. Da das Herauslesen aufgabenspezifischer Episode-Level-Daten zu Zielerreichung und Kollision aus dem aufgezeichneten Datenbestand im Rahmen dieser Machbarkeitsstudie nicht erfolgreich umgesetzt werden konnte (vgl. \autoref{ch:results}), stützt sich die qualitative Beurteilung der Zielerreichung, Zeit bis zum Ziel, Kollisionsvermeidung und Zieldistanz auf die visuelle Analyse der Aufnahmen. Diese qualitativen Befunde werden durch die folgenden quantitativen Trainingsmetriken untermauert:
\begin{itemize}
	\item \textbf{Kumulierter Return ($R$):} Die über die Episode aufsummierte Gesamtrendite. Der Return dient als Optimierungsziel für den PPO-Algorithmus, wird jedoch \textbf{nicht als alleinige Erfolgsmetrik} herangezogen, da sein absoluter Betrag stark von der gewählten Skalierung der Teilbelohnungen und Strafterme abhängt und nicht direkt proportional zur Aufgabenbewältigung ist.
	\item \textbf{Mittlere Episodenlänge ($\bar{L}_{\text{ep}}$):} Die durchschnittliche Dauer einer Episode. Ein plötzlicher Einbruch dieser Metrik während des Trainings ist ein kritischer Indikator für Instabilität oder vorzeitige Episodenabbrüche durch Stürze (z.\,B. Überschreiten kritischer Roll- oder Nickwinkel).
	\item \textbf{Policy-Entropie ($H(\pi)$):} Die Entropie der Aktionsverteilung liefert diagnostische Einblicke in die Balance zwischen Exploration (hohe Entropie) und deterministischer Verfestigung des Verhaltens (sinkende Entropie).
	\item \textbf{Curriculumspezifische Belohnungsterme ($r_{\text{term}}$):} Die für die jeweilige Curriculum-Stufe konfigurierten Teilbelohnungen (z.\,B. Zielannäherung, Sprungphasen, Stabilität) dienen der Zuordnung beobachteter Verhaltensweisen zu den intendierten Lernzielen und werden stets in Relation zu Return und Episodenlänge betrachtet.
\end{itemize}

Für Trainingskurven mit mehreren vorliegenden Läufen werden Mittelwert und Streubreite über die dokumentierten Seeds ausgewertet. Einzelne Kurven oder Checkpoints werden als Einzelbeobachtungen gekennzeichnet. Die qualitative Inspektion ausgewählter Episoden anhand von Videoaufzeichnungen bildet die Grundlage der aufgabenspezifischen Verhaltensbeurteilung und wird durch die genannten Trainingsmetriken untermauert; dabei werden Bewegungsmuster, Schrittkoordination und das Verhalten bei Hinderniskontakten im Detail analysiert. Die aufgezeichneten Videos sind online dokumentiert.\footnote{Videodokumentation der experimentellen Ergebnisse: \url{https://drive.google.com/drive/folders/1ZInYQzLSj-5zsYY7pxdSiBy7aekxNdNQ?usp=drive_link}}